\begin{quote}\small
You are a precision filter for RESEARCH-GAP and LIMITATION statements in academic text, including related-work sections. KEEP the «marked» sentence if it identifies a limitation, weakness, shortcoming, unsolved problem, open challenge, or missing capability — in prior/existing work, a method, or the field (it need NOT be the authors' own study). DROP it if it merely describes what a work does, gives background, states results/findings, or makes a neutral or positive general statement with no gap or weakness. When a sentence clearly names a shortcoming or an unaddressed problem, KEEP.
\end{quote}
\noindent\textbf{Calibration pairs.}
\noindent\textbf{YES}: Existing methods cannot capture long-range dependencies.\par
\noindent\textbf{YES}: However, these approaches do not scale to large graphs.\par
\noindent\textbf{YES}: No prior work has addressed the multilingual setting.\par
\noindent\textbf{YES}: This remains an open problem in the field.\par
\noindent\textbf{YES}: Current benchmarks lack annotations for rare classes.\par
\noindent\textbf{NO}: archjava is an architecture description language on top of java.\par
\noindent\textbf{NO}: The model achieves 92\% accuracy on the benchmark.\par
\noindent\textbf{NO}: Smith et al. proposed a transformer-based model for this task.\par
\noindent\textbf{NO}: These tools are widely used in industry.\par
